\section{The class \texorpdfstring{\(\PSH_X(\theta)\)}{PSH(theta)}}
\label{sec:ch34-psh}

Everything so far has been finitary: a Fubini--Study function is built from
finitely many sections, a PL function from a single test configuration.  The
class of psh functions is obtained by closing up under decreasing limits,
and it is the smallest class for which this is done.

\begin{definition}
\label{def:bj-psh}
Let \(\theta\in N^1(X)\).  A function
\(\phi:X^{\An}\to\mathbf R\cup\{-\infty\}\) is \emph{\(\theta\)-psh} if
\begin{enumerate}
\item \(\phi\) is upper semicontinuous;
\item \(\phi\) is generically finite, that is, \(\phi\not\equiv-\infty\) on
      any irreducible component of \(X\); and
\item there are \(\mathbf Q\)-line bundles \(L_i\) and a decreasing net
      \(\phi_i\in\cH^{\mathrm{gf}}_{\mathbf R}(L_i)\) with
      \(\lim_i c_1(L_i)=\theta\) in \(N^1(X)\) and \(\phi=\lim_i\phi_i\)
      pointwise on \(X^{\An}\).
\end{enumerate}
The set of \(\theta\)-psh functions is denoted \(\PSH_X(\theta)\).
\end{definition}

Three points of the definition deserve comment before anything is proved
with it.

First, condition~(3) uses a \emph{net}, not a sequence.  This is not
pedantry: \(X^{\An}\) is compact but not metrizable in general, the index
set of the natural approximating family is the set of test configurations
ordered by domination, and that set is directed but not countable.  The
places where a sequence \emph{may} legitimately be extracted are recorded in
\Cref{chap:bj-linear-growth}; none of them is available here.

Second, the limit in~(3) is a decreasing limit, so it exists in
\(\Rbar\) at every point without hypothesis, by ledger entry~(1), and the
limit function is automatically upper semicontinuous as an infimum of
continuous functions.  Condition~(1) is therefore redundant given~(3); it is
retained because it is the condition one verifies in practice when a
candidate \(\phi\) is presented by other means.

Third, condition~(2) is a genuine nontriviality condition.  It excludes the
function \(\equiv-\infty\), and on a reducible \(X\) it is imposed
componentwise: being finite on only some components is not enough.

Fourth, and most importantly, the classes are allowed to \emph{vary}: the
approximants live in \(\cH^{\mathrm{gf}}_{\mathbf R}(L_i)\) for line bundles
whose classes merely converge to \(\theta\).  This is not a technical
convenience: it is part of the closure construction of
\cite[Definition~4.1]{BJ22} and allows arbitrary numerical limit classes.
No claim is made that fixed-class approximation would force nefness; GF
Fubini--Study functions can occur in non-nef effective classes.  For \(E\)
an effective divisor, \(\log|s_E|\) is \([E]\)-psh, and \([E]\) is only
pseudoeffective in general.  The precise
consequence of the varying-class convention is
\Cref{lem:ch34-psh-nonempty}, which separates the pseudoeffective threshold
for \(\PSH_X(\theta)\) from the nef threshold for its \emph{bounded}
members.

\begin{lemma}
\label{lem:ch34-psh-nonempty}
Let \(\theta\in N^1(X)\).  Then
\begin{equation}
\label{eq:psh-existence-cones}
 \begin{aligned}
 \PSH_X(\theta)\ne\emptyset
 &\implies \theta\in\operatorname{Psef}(X),\\
 \PSH_X(\theta)\ \text{contains a bounded function}
 &\implies \theta\in\Nef(X),
 \end{aligned}
\end{equation}
and conversely the constants are \(\theta\)-psh whenever
\(\theta\in\Nef(X)\).  In particular
\(\mathbf R\subset\PSH_X(\theta)\) if and only if \(\theta\in\Nef(X)\).
\end{lemma}

\begin{proof}
Suppose \(\PSH_X(\theta)\ne\emptyset\) and let \((\phi_i)\) be as in
\Cref{def:bj-psh}(3).  Each \(\cH^{\mathrm{gf}}_{\mathbf R}(L_i)\) is
nonempty, so \(L_i\) admits a nonzero section after clearing denominators
and is therefore effective; hence \(c_1(L_i)\in\operatorname{Psef}(X)\).  The
pseudoeffective cone is closed, and \(c_1(L_i)\to\theta\), so
\(\theta\in\operatorname{Psef}(X)\).

Suppose now that some \(\phi\in\PSH_X(\theta)\) is bounded.  Every member
\(\phi_i\) of a defining decreasing net satisfies \(\phi_i\geq\phi\), so it
is finite valued, hence
\(\phi_i\in\cH_{\mathbf R}(L_i)\) and the defining sections have no common
zero; consequently \(L_i\) is semiample, so \(c_1(L_i)\in\Nef(X)\).  The nef
cone is closed, whence \(\theta\in\Nef(X)\).

Conversely let \(\theta\in\Nef(X)\) and \(c\in\mathbf R\).  Choose
\(\mathbf Q\)-line bundles \(L_i\) with \(c_1(L_i)\to\theta\) and \(L_i\)
ample; for each \(i\) the base-point-free sections of a large multiple
present the constant \(c\) as an element of \(\cH_{\mathbf R}(L_i)\) after
translating the weights, and the net may be arranged to decrease to \(c\).
Hence \(c\in\PSH_X(\theta)\).  Since constants are bounded, the second
implication of \eqref{eq:psh-existence-cones} shows the converse fails when
\(\theta\) is not nef.
\end{proof}
